% Generated draft: requires mathematical and source review.
\documentclass[11pt]{article}
\usepackage[margin=25mm]{geometry}
\usepackage{fontspec}
\setmainfont{Latin Modern Roman}
\usepackage{amsmath,amssymb,mathtools,unicode-math}
\setmathfont{Latin Modern Math}
\usepackage{enumitem,needspace,xurl,hyperref}
\setlength{\parindent}{0pt}
\setlength{\parskip}{5pt}
\setlength{\emergencystretch}{3em}
\newcommand{\N}{\mathbb{N}}
\newcommand{\Z}{\mathbb{Z}}
\newcommand{\Q}{\mathbb{Q}}
\newcommand{\R}{\mathbb{R}}
\newcommand{\C}{\mathbb{C}}
\newcommand{\F}{\mathbb{F}}
\newcommand{\A}{\mathbb{A}}
\newcommand{\PP}{\mathbb{P}}
\newcommand{\E}{\mathbb{E}}
\newcommand{\eps}{\varepsilon}
\newcommand{\dd}{\,\mathrm{d}}
\begin{document}
\section*{UnsolvedMath ID: 2100205}
\textbf{Source problem code:} catalogue entry
\subsection{Definitions and mathematical statement}
Consider a Liouville-integrable Hamiltonian system on a symplectic $2n$-manifold $(M,\omega)$, with commuting first integrals $F=(F_1,\ldots,F_n)$ independent on the regular locus. Near compact regular Lagrangian torus fibers, action variables $I_1,\ldots,I_n$ are locally defined by periods
\[
I_j(b)=\frac{1}{2\pi}\int_{\gamma_j(b)}\lambda,
\]
where $\dd\lambda=\omega$ locally and the $\gamma_j(b)$ form an integral cycle basis on the torus over a regular value $b$. This is the usual local action convention; no global single-valued action coordinates or global primitive is assumed. Cycle changes alter the actions by integral linear changes, and normalization permits additive constants; branches and monodromy may matter near a singular value.

Let $c$ be a singular value and $L_c=F^{-1}(c)$ the singular fiber under consideration. Assume, as in the source, that explicit formulas for the actions are available and their asymptotics can be analyzed as regular values approach $c$. Can these asymptotics recover the topology of the singularity, or at least distinguish its types? A precise criterion would have to specify which asymptotic data are retained and which topology is to be recovered: the fiber itself, a germ of the singular Lagrangian fibration, or a neighborhood of the entire singular fiber. It would also have to account for permissible action and base-coordinate changes. The record does not fix those equivalence conventions or restrict to nondegenerate singularities in the general question.

The source asks specifically whether the action singular part displayed schematically as
\[
I_{\mathrm{sing}}(h)=h\ln h+\cdots
\]
is characteristic of a nondegenerate hyperbolic singularity. For a real logarithm this displayed expression is on a branch $h>0$; the record does not define the transverse parameter $h$, the normalization of the coefficient, the omitted terms, or how other variables are held fixed. A two-sided formulation would require an explicit convention such as $\ln|h|$, which is not silently substituted. The converse question is whether the appropriate invariantly formulated asymptotic signature forces the intended nondegenerate hyperbolic type, not merely whether hyperbolic examples have logarithmic terms.

The supplied summary reports classical logarithmic hyperbolic asymptotics and partial distinctions among elliptic, hyperbolic, and focus--focus cases, but gives no complete all-types criterion. Its assertions do not imply that an unspecified logarithmic term alone separates all these types. The original says the one-degree-of-freedom case is understood but has a missing citation, ending ``in .''. These are attributed partial-progress claims, not a classification or proof supplied here.
\subsection{Short English statement}
Can the asymptotic behavior of explicitly known action variables identify a singular integrable fiber or its topological type? In particular, does the correctly normalized hyperbolic action signature characterize nondegenerate hyperbolic singularities? The source gives a schematic logarithmic term but omits the asymptotic equivalence, parameter conventions, and the citation for its claimed one-degree-of-freedom result.
\subsection{Sources}
\begin{itemize}
\item Local source registry: \url{data/unsolvedmath/problems.json}, permanent UnsolvedMath ID $2100205$, source code catalogue entry, including the supplied PARTIAL-PROGRESS summary. Dataset attribution: Ulam AI and UnsolvedMath Contributors; original-source attribution and rights are retained.
\item \emph{Open Problems, Questions, and Challenges in Finite-Dimensional Integrable Systems} (2018), Question 2.5: \url{https://arxiv.org/abs/1804.03737}. The supplied metadata reports source-TeX extraction and July 29, 2026 access, with current-status and rights review flags. No author list is supplied in this record.
\item The summary mentions classical ``hyperbolic monodromy'' work and Zung's and others' analyses of singular fibers, but gives no precise additional bibliography. The original source's reference for the one-degree-of-freedom case is missing from the extract. No linked source was inspected for this draft.
\end{itemize}
\end{document}
